% !PRINTED 308
\weberpage{308}
%<<CONT
\[
\left|\begin{array}{c|cc|cc|cc|cc}
x-0 & PV, & PV & PP, & PV & VP, & VP & VV, & VP \\ \hline
x+0 & PP, & PP & PV, & PP & VV, & VV & VP, & VV
\end{array}\right|
\text{.}
\]
Im ersten Falle werden zwei $PP$ \gesperrt{gewonnen}, im zweiten und vierten Falle tritt keine Aenderung ein, im dritten Falle gehen zwei $VP$ \gesperrt{verloren}.

Damit aber sind die Sätze II, III bewiesen.

Die Voraussetzungen 2., 3., dass von den Functionen $f_\nu$ und $F_\nu$ nicht zwei benachbarte für dasselbe $x$ verschwinden sollen, können wir aufgeben. Denn wenn wir an der ersten Voraussetzung festhalten, dass keine mehrfachen Wurzeln vorhanden sind, so können wir dem Coëfficienten $a_\nu$ so kleine Aenderungen ertheilen, dass erstens die Voraussetzungen 2., 3. befriedigt sind, und zweitens die zwischen $\alpha$ und $\beta$ gelegenen Wurzeln sich so wenig ändern, dass keine von ihnen aus dem Intervall heraustritt. Da die Wurzeln nur einzeln vorkommen, so können dabei auch nicht reelle Wurzeln in imaginäre übergehen.

Ob der Satz bei richtiger Zählung der mehrfachen Wurzeln auch noch im Falle mehrfacher Wurzeln gültig bleibt, mag dahin gestellt bleiben.

\begin{center}
\S.\,101.

Der Cartesische Lehrsatz.
\end{center}

Der \gesperrt{Cartesische} Lehrsatz, der auch nach \gesperrt{Harriot}\footnote{Harriot, Artis analyticae praxis ad aequationes algebraicas resolvendas 1631. Gauss' Werke, Bd. III, S. 67. Lagrange, l. c.} benannt wird, giebt eine Regel zur Abschätzung der Zahl der positiven Wurzeln. Wir können ihn einfach als speciellen Fall des \gesperrt{Budan}'schen Theorems betrachten, indem wir $\alpha = 0$ und $\beta = \infty$ oder wenigstens so gross annehmen, dass die Functionen
\weq{1}{f(x), f'(x), f''(x) \ldots f^{(n)}(x)}
für $x > \beta$ alle dasselbe Zeichen haben. Dies Zeichen stimmt überein mit dem Zeichen des Coëfficienten der höchsten Potenz von $x$ in $f(x)$ und kann positiv angenommen werden.

Ist
\[
f(x) = a_0 x^n + a_1 x^{n-1} + a_2 x^{n-2} + \cdots + a_{n-1} x + a_n,
\]
%CONT>>
